\section{制度与环境}

西周不是后来那种“中央下令、地方照办”的帝国。周天子靠宗族关系、封国、婚姻、赏赐和出兵，把许多原本独立的政治共同体绑在一起。王畿、同姓诸侯、功臣和旧商贵族既是一张网里的伙伴，也随时可能成为对手；青铜礼器和册命文书不只是好看的礼物，更是在确认谁有权、谁有功、谁该服从。

王室的强处，是能把核心地区资源、宗族合法性和军事动员拉到一起；它的弱处也在这里：封国一旦坐大，王畿资源一旦紧张，或者继承出了乱子，周天子没有后世郡县国家那样直接管到地方的手。

\section{先用一条线看完西周}
\begin{center}
\begin{tikzpicture}[x=.027cm,y=1cm]
  \draw[line width=.8pt,color=timeline] (0,0) -- (275,0);
  \draw[color=dynasty,line width=.8pt] (0,-.11) -- (0,.11);
  \draw[color=dynasty,line width=.8pt] (69,-.11) -- (69,.11);
  \draw[color=dynasty,line width=.8pt] (205,-.11) -- (205,.11);
  \draw[color=dynasty,line width=.8pt] (275,-.11) -- (275,.11);
  \node[above,align=center,font=\small\color{ink}] at (0,0) {约前1046\\牧野};
  \node[below,align=center,font=\small\color{ink}] at (69,0) {前977\\昭王南征};
  \node[above,align=center,font=\small\color{ink}] at (205,0) {前841\\共和};
  \node[below,align=center,font=\small\color{ink}] at (275,0) {前771\\镐京陷落};
\end{tikzpicture}
\end{center}

\begin{sourcenote}[这一章主要参考什么]
西周叙事以\sourcebook{史记·周本纪}、\sourcebook{尚书}、\sourcebook{国语}和\sourcebook{竹书纪年}等传世材料为入口，并须参考金文与现代断代研究。越早的事件，越不能把后来的整齐叙事当成毫无疑问的现场记录。
\end{sourcenote}
